\chapter{Detailed Off-Chain Timing Tables}
\label{app:offchain-timing-tables}

This appendix gathers the off-chain timings that are used in the discussion of large-file
execution. They are wall-clock timings on the benchmark machine, not gas measurements. The
values are therefore useful for relative interpretation and reproducibility, but they are
not protocol constants.

\section{Native 1 GiB benchmark timings}

\begin{table}[H]
  \centering
  \begin{tabularx}{\linewidth}{>{\raggedright\arraybackslash}p{0.30\linewidth}r>{\raggedright\arraybackslash}X}
    \toprule
    Measurement & Time & Interpretation \\
    \midrule
    Precontract CLI & 22.42--48.70 s & Vendor-side native preparation of ciphertext,
    commitments, circuit metadata, and opening values. \\
    Native dispute generation & 16.43--35.46 s & Aggregate native benchmark time for preparing
    dispute data; not attributable to one actor only. \\
    Total \(hpre_i\) generation & 9.47--11.92 s & Buyer-side generation of challenge responses
    over 26 rounds in the forced-left benchmark. \\
    Local Hardhat transactions & 0.30--0.56 s & Local transaction submission time in the test
    environment; not a public-chain latency measurement. \\
    \bottomrule
  \end{tabularx}
  \caption{Observed off-chain timing ranges for repeated 1 GiB hardcoded benchmark runs.}
  \label{tab:app-1gib-timings}
\end{table}

\section[First hpre rounds]{First forced-left \(hpre_i\) rounds}

\begin{table}[H]
  \centering
  \begin{tabular}{rrr}
    \toprule
    Round & Challenge index \(i\) & Native \(hpre_i\) time \\
    \midrule
    0 & 16,777,219 & 4.58--6.04 s \\
    1 & 8,388,610 & 2.23--2.46 s \\
    2 & 4,194,305 & 1.17--1.33 s \\
    3 & 2,097,153 & 0.60--0.71 s \\
    4 & 1,048,577 & 0.31--0.39 s \\
    \bottomrule
  \end{tabular}
  \caption{Observed first \(hpre_i\) timing ranges in repeated 1 GiB forced-left dispute benchmarks.}
  \label{tab:app-hpre-rounds}
\end{table}

The decreasing times in Table~\ref{tab:app-hpre-rounds} are CPU times, not gas. The
benchmark forces the left branch of the binary search, so the challenge index is
approximately divided by two at each round. The native computation recomputes a smaller
prefix for later rounds.

\section{Native versus browser/WASM interpretation}

\begin{table}[H]
  \centering
  \begin{tabularx}{\linewidth}{>{\raggedright\arraybackslash}p{0.24\linewidth}>{\raggedright\arraybackslash}X}
    \toprule
    Path & Role in the project \\
    \midrule
    WebAssembly frontend path & Useful for integration, UI demonstrations, and small or
    medium files. It is not the reliable path for 1 GiB measurements. \\
    Native Rust CLI path & Used for large-file precontract generation and hardcoded
    dispute witnesses. The 1 GiB results in this report are native CLI results. \\
    Hardhat transaction time & Measures local execution of test transactions. It does not
    include public-network latency, mempool behavior, or wallet interaction delays. \\
    \bottomrule
  \end{tabularx}
  \caption{Interpretation of the off-chain execution paths.}
  \label{tab:app-native-vs-wasm}
\end{table}

In the direct WASM experiments, files around 640 MiB and 704 MiB could be processed after
memory-related patches, while larger sizes such as 736 MiB remained unreliable in that path.
For this reason, the report claims 1 GiB support only for the native CLI benchmark path.
